% ============================================================
% Question Bank: ch07-01 Parts of a Circle
% Topic Title: Parts of a Circle
% CCSS: 7.G.B.4
% Questions: 30 (18 MC + 7 SA + 5 Graphical)
% ============================================================

% --- Multiple Choice Questions (18) ---

\begin{practiceQuestion}{7-1-q01}{mc}
\begin{questionText}
What is the name of the distance from the center of a circle to any point on the circle?
\end{questionText}
\choiceA{Diameter}
\choiceB{Circumference}
\choiceC{Radius}
\choiceD{Chord}
\correctAnswer{C}
\explanation{The radius is the distance from the center of a circle to any point on the circle.}
\end{practiceQuestion}

\begin{practiceQuestion}{7-1-q02}{mc}
\begin{questionText}
A circle has a radius of $6$ cm. What is its diameter?
\end{questionText}
\choiceA{$3$ cm}
\choiceB{$6$ cm}
\choiceC{$12$ cm}
\choiceD{$18$ cm}
\correctAnswer{C}
\explanation{The diameter is twice the radius: $d = 2r = 2 \times 6 = 12$ cm.}
\end{practiceQuestion}

\begin{practiceQuestion}{7-1-q03}{mc}
\begin{questionText}
A circle has a diameter of $14$ inches. What is its radius?
\end{questionText}
\choiceA{$28$ inches}
\choiceB{$14$ inches}
\choiceC{$7$ inches}
\choiceD{$3.5$ inches}
\correctAnswer{C}
\explanation{The radius is half the diameter: $r = d \div 2 = 14 \div 2 = 7$ inches.}
\end{practiceQuestion}

\begin{practiceQuestion}{7-1-q04}{mc}
\begin{questionText}
Which of the following best describes the circumference of a circle?
\end{questionText}
\choiceA{The distance from the center to the edge}
\choiceB{The distance across the circle through the center}
\choiceC{The distance around the circle}
\choiceD{The area inside the circle}
\correctAnswer{C}
\explanation{The circumference is the distance around (perimeter of) a circle.}
\end{practiceQuestion}

\begin{practiceQuestion}{7-1-q05}{mc}
\begin{questionText}
A circle has a diameter of $9$ feet. What is its radius?
\end{questionText}
\choiceA{$18$ feet}
\choiceB{$9$ feet}
\choiceC{$4.5$ feet}
\choiceD{$3$ feet}
\correctAnswer{C}
\explanation{The radius is half the diameter: $r = 9 \div 2 = 4.5$ feet.}
\end{practiceQuestion}

\begin{practiceQuestion}{7-1-q06}{mc}
\begin{questionText}
Which statement about the diameter and radius of a circle is always true?
\end{questionText}
\choiceA{The diameter is half the radius}
\choiceB{The diameter equals the radius}
\choiceC{The diameter is twice the radius}
\choiceD{The diameter is three times the radius}
\correctAnswer{C}
\explanation{The diameter is always twice the radius: $d = 2r$.}
\end{practiceQuestion}

\begin{practiceQuestion}{7-1-q07}{mc}
\begin{questionText}
A circle has a radius of $3.5$ meters. What is its diameter?
\end{questionText}
\choiceA{$1.75$ m}
\choiceB{$3.5$ m}
\choiceC{$7$ m}
\choiceD{$10.5$ m}
\correctAnswer{C}
\explanation{$d = 2r = 2 \times 3.5 = 7$ m.}
\end{practiceQuestion}

\begin{practiceQuestion}{7-1-q08}{mc}
\begin{questionText}
The diameter of a bicycle wheel is $26$ inches. What is the radius of the wheel?
\end{questionText}
\choiceA{$52$ inches}
\choiceB{$26$ inches}
\choiceC{$13$ inches}
\choiceD{$6.5$ inches}
\correctAnswer{C}
\explanation{$r = d \div 2 = 26 \div 2 = 13$ inches.}
\end{practiceQuestion}

\begin{practiceQuestion}{7-1-q09}{mc}
\begin{questionText}
A line segment that passes through the center of a circle and has both endpoints on the circle is called a \underline{\hspace{1.5cm}}.
\end{questionText}
\choiceA{Radius}
\choiceB{Chord}
\choiceC{Diameter}
\choiceD{Circumference}
\correctAnswer{C}
\explanation{A diameter passes through the center and connects two points on the circle.}
\end{practiceQuestion}

\begin{practiceQuestion}{7-1-q10}{mc}
\begin{questionText}
If the radius of a circle is $\frac{3}{4}$ inch, what is the diameter?
\end{questionText}
\choiceA{$\frac{3}{8}$ inch}
\choiceB{$\frac{3}{4}$ inch}
\choiceC{$1\frac{1}{2}$ inches}
\choiceD{$3$ inches}
\correctAnswer{C}
\explanation{$d = 2r = 2 \times \frac{3}{4} = \frac{6}{4} = 1\frac{1}{2}$ inches.}
\end{practiceQuestion}

\begin{practiceQuestion}{7-1-q11}{mc}
\begin{questionText}
Which part of a circle is labeled $O$ in most diagrams?
\end{questionText}
\choiceA{A point on the circumference}
\choiceB{The center}
\choiceC{The diameter}
\choiceD{The radius}
\correctAnswer{B}
\explanation{In standard circle diagrams, $O$ typically labels the center of the circle.}
\end{practiceQuestion}

\begin{practiceQuestion}{7-1-q12}{mc}
\begin{questionText}
A clock face is a circle with a radius of $15$ cm. What is the diameter of the clock face?
\end{questionText}
\choiceA{$7.5$ cm}
\choiceB{$15$ cm}
\choiceC{$30$ cm}
\choiceD{$45$ cm}
\correctAnswer{C}
\explanation{$d = 2r = 2 \times 15 = 30$ cm.}
\end{practiceQuestion}

\begin{practiceQuestion}{7-1-q13}{mc}
\begin{questionText}
A circle has a circumference of $31.4$ cm. Using $\pi \approx 3.14$, what is the diameter?
\end{questionText}
\choiceA{$5$ cm}
\choiceB{$10$ cm}
\choiceC{$15$ cm}
\choiceD{$20$ cm}
\correctAnswer{B}
\explanation{$d = C \div \pi = 31.4 \div 3.14 = 10$ cm.}
\end{practiceQuestion}

\begin{practiceQuestion}{7-1-q14}{mc}
\begin{questionText}
How many radii does it take to make one diameter?
\end{questionText}
\choiceA{$1$}
\choiceB{$2$}
\choiceC{$3$}
\choiceD{$4$}
\correctAnswer{B}
\explanation{A diameter is made up of two radii placed end to end through the center.}
\end{practiceQuestion}

\begin{practiceQuestion}{7-1-q15}{mc}
\begin{questionText}
The diameter of a pizza is $16$ inches. What is the radius?
\end{questionText}
\choiceA{$4$ inches}
\choiceB{$8$ inches}
\choiceC{$16$ inches}
\choiceD{$32$ inches}
\correctAnswer{B}
\explanation{$r = d \div 2 = 16 \div 2 = 8$ inches.}
\end{practiceQuestion}

\begin{practiceQuestion}{7-1-q16}{mc}
\begin{questionText}
A circle has a diameter of $2.4$ m. What is the radius?
\end{questionText}
\choiceA{$0.6$ m}
\choiceB{$1.2$ m}
\choiceC{$2.4$ m}
\choiceD{$4.8$ m}
\correctAnswer{B}
\explanation{$r = d \div 2 = 2.4 \div 2 = 1.2$ m.}
\end{practiceQuestion}

\begin{practiceQuestion}{7-1-q17}{mc}
\begin{questionText}
Which of the following is NOT a part of a circle?
\end{questionText}
\choiceA{Center}
\choiceB{Radius}
\choiceC{Vertex}
\choiceD{Diameter}
\correctAnswer{C}
\explanation{A vertex is a corner point found on polygons, not circles. Center, radius, and diameter are all parts of a circle.}
\end{practiceQuestion}

\begin{practiceQuestion}{7-1-q18}{mc}
\begin{questionText}
A circular garden has a radius of $8$ feet. A path goes straight through the center from one side to the other. How long is the path?
\end{questionText}
\choiceA{$4$ feet}
\choiceB{$8$ feet}
\choiceC{$16$ feet}
\choiceD{$24$ feet}
\correctAnswer{C}
\explanation{The path through the center is a diameter: $d = 2r = 2 \times 8 = 16$ feet.}
\end{practiceQuestion}

% --- Short Answer Questions (7) ---

\begin{practiceQuestion}{7-1-q19}{sa}
\begin{questionText}
A circle has a radius of $11$ cm. What is its diameter?
\end{questionText}
\correctAnswer{$22$ cm}
\explanation{$d = 2r = 2 \times 11 = 22$ cm.}
\end{practiceQuestion}

\begin{practiceQuestion}{7-1-q20}{sa}
\begin{questionText}
A circle has a diameter of $5.6$ inches. What is the radius?
\end{questionText}
\correctAnswer{$2.8$ inches}
\explanation{$r = d \div 2 = 5.6 \div 2 = 2.8$ inches.}
\end{practiceQuestion}

\begin{practiceQuestion}{7-1-q21}{sa}
\begin{questionText}
A circular table has a diameter of $4$ feet. What is the radius of the table in inches? (Hint: $1$ foot $= 12$ inches)
\end{questionText}
\correctAnswer{$24$ inches}
\explanation{Diameter $= 4$ feet $= 48$ inches. Radius $= 48 \div 2 = 24$ inches.}
\end{practiceQuestion}

\begin{practiceQuestion}{7-1-q22}{sa}
\begin{questionText}
The circumference of a circle is $62.8$ cm. Using $\pi \approx 3.14$, find the radius.
\end{questionText}
\correctAnswer{$10$ cm}
\explanation{$d = C \div \pi = 62.8 \div 3.14 = 20$ cm. Then $r = 20 \div 2 = 10$ cm.}
\end{practiceQuestion}

\begin{practiceQuestion}{7-1-q23}{sa}
\begin{questionText}
A circle has a radius of $\frac{5}{2}$ meters. What is the diameter?
\end{questionText}
\correctAnswer{$5$ meters}
\explanation{$d = 2r = 2 \times \frac{5}{2} = 5$ meters.}
\end{practiceQuestion}

\begin{practiceQuestion}{7-1-q24}{sa}
\begin{questionText}
Two circles have radii of $4$ cm and $9$ cm. What is the difference between their diameters?
\end{questionText}
\correctAnswer{$10$ cm}
\explanation{Diameters: $2 \times 4 = 8$ cm and $2 \times 9 = 18$ cm. Difference: $18 - 8 = 10$ cm.}
\end{practiceQuestion}

\begin{practiceQuestion}{7-1-q25}{sa}
\begin{questionText}
A circle has a circumference of $47.1$ cm. Using $\pi \approx 3.14$, what is the diameter?
\end{questionText}
\correctAnswer{$15$ cm}
\explanation{$d = C \div \pi = 47.1 \div 3.14 = 15$ cm.}
\end{practiceQuestion}

% --- Graphical Questions (3 GMC + 2 GSA) ---

\begin{practiceQuestion}{7-1-q26}{gmc}
\begin{questionText}
Look at the circle below. Which line segment represents the radius?

\begin{center}
\begin{tikzpicture}[scale=0.8]
  \draw[thick] (0,0) circle (2);
  \fill (0,0) circle (2pt) node[below left] {$O$};
  \draw[thick,funBlue] (0,0) -- (2,0) node[midway,above] {$\overline{OA}$};
  \fill[funBlue] (2,0) circle (2pt) node[right] {$A$};
  \draw[thick,funRed] (-2,0) -- (2,0) node[midway,below=4pt] {$\overline{BA}$};
  \fill[funRed] (-2,0) circle (2pt) node[left] {$B$};
  \draw[thick,funGreen,dashed] (0,0) -- (0,2);
  \fill[funGreen] (0,2) circle (2pt) node[above] {$C$};
  \node[right=4pt,funGreen] at (0,1) {$\overline{OC}$};
\end{tikzpicture}
\end{center}
\end{questionText}
\choiceA{$\overline{BA}$ only}
\choiceB{$\overline{OA}$ and $\overline{OC}$}
\choiceC{$\overline{BA}$ and $\overline{OC}$}
\choiceD{$\overline{OA}$ only}
\correctAnswer{B}
\explanation{Both $\overline{OA}$ and $\overline{OC}$ go from the center $O$ to a point on the circle, so both are radii. $\overline{BA}$ passes through the center and is a diameter.}
\end{practiceQuestion}

\begin{practiceQuestion}{7-1-q27}{gmc}
\begin{questionText}
Use the diagram below. What is the diameter of the circle?

\begin{center}
\begin{tikzpicture}[scale=0.7]
  \draw[thick] (0,0) circle (2.5);
  \fill (0,0) circle (2pt) node[below left] {$O$};
  \draw[thick,funBlue,<->] (-2.5,0) -- (2.5,0);
  \node[below] at (-2.5,0) {$P$};
  \node[below] at (2.5,0) {$Q$};
  \draw[thick,funRed,<->] (0,0) -- (0,2.5);
  \node[left,funRed] at (0,1.25) {$7$ cm};
\end{tikzpicture}
\end{center}
\end{questionText}
\choiceA{$3.5$ cm}
\choiceB{$7$ cm}
\choiceC{$14$ cm}
\choiceD{$21$ cm}
\correctAnswer{C}
\explanation{The red segment from $O$ to the top is a radius of $7$ cm. The diameter is $2 \times 7 = 14$ cm.}
\end{practiceQuestion}

\begin{practiceQuestion}{7-1-q28}{gmc}
\begin{questionText}
In the diagram, circle $O$ has a diameter of $20$ cm. What is the length of $\overline{OM}$?

\begin{center}
\begin{tikzpicture}[scale=0.6]
  \draw[thick] (0,0) circle (3);
  \fill (0,0) circle (2pt) node[below] {$O$};
  \draw[thick,funBlue] (0,0) -- (3,0);
  \fill[funBlue] (3,0) circle (2pt) node[right] {$M$};
  \draw[thick,funRed,<->] (-3,-0.6) -- (3,-0.6);
  \node[below,funRed] at (0,-0.6) {$20$ cm};
\end{tikzpicture}
\end{center}
\end{questionText}
\choiceA{$5$ cm}
\choiceB{$10$ cm}
\choiceC{$20$ cm}
\choiceD{$40$ cm}
\correctAnswer{B}
\explanation{$\overline{OM}$ is a radius. Radius $= \text{diameter} \div 2 = 20 \div 2 = 10$ cm.}
\end{practiceQuestion}

\begin{practiceQuestion}{7-1-q29}{gsa}
\begin{questionText}
Look at the two concentric circles (circles with the same center). The radius of the smaller circle is $3$ cm and the radius of the larger circle is $5$ cm. What is the difference between their diameters?

\begin{center}
\begin{tikzpicture}[scale=0.6]
  \draw[thick,funBlue] (0,0) circle (1.5);
  \draw[thick,funRed] (0,0) circle (2.5);
  \fill (0,0) circle (2pt) node[below left] {$O$};
  \draw[thick,<->] (0,0) -- (1.5,0) node[midway,above] {$3$};
  \draw[thick,<->] (0,0) -- (-2.5,0) node[midway,above] {$5$};
  \node[below right] at (1.5,0) {cm};
  \node[below left] at (-2.5,0) {cm};
\end{tikzpicture}
\end{center}
\end{questionText}
\correctAnswer{$4$ cm}
\explanation{Smaller diameter: $2 \times 3 = 6$ cm. Larger diameter: $2 \times 5 = 10$ cm. Difference: $10 - 6 = 4$ cm.}
\end{practiceQuestion}

\begin{practiceQuestion}{7-1-q30}{gsa}
\begin{questionText}
A circular running track is shown below. The diameter across the track is $50$ meters. Using $\pi \approx 3.14$, what is the circumference of the track?

\begin{center}
\begin{tikzpicture}[scale=0.5]
  \draw[thick,funBlue,fill=funBlue!10] (0,0) circle (3);
  \draw[thick,funRed,<->] (-3,0) -- (3,0);
  \node[above,funRed] at (0,0) {$50$ m};
  \fill (0,0) circle (2pt);
\end{tikzpicture}
\end{center}
\end{questionText}
\correctAnswer{$157$ m}
\explanation{$C = \pi d = 3.14 \times 50 = 157$ m.}
\end{practiceQuestion}
